\documentclass{article}
\usepackage{graphicx} % Required for inserting images
\usepackage{xcolor}

\usepackage{amsmath}
\usepackage{amsthm}
\usepackage{thmtools} 
\usepackage{cleveref}


% Number equations within sections (instead of chapters)
\numberwithin{equation}{section}

% Declare theorem environments anchored to sections
\declaretheorem[name=Theorem,numberwithin=section]{theorem}
\declaretheorem[name=Lemma,numberlike=theorem]{lemma}
\declaretheorem[name=Proposition,numberlike=theorem]{proposition}
\declaretheorem[name=Corollary,numberlike=theorem]{corollary}
\declaretheorem[name=Definition,numberlike=theorem]{definition}
\declaretheorem[name=Assumption,numberlike=theorem]{assumption}
\declaretheorem[name=Example,numberlike=theorem]{example}
\declaretheorem[name=Remark,numberlike=theorem]{remark}

% For restatable theorems 
\declaretheorem[name=Theorem,numberwithin=section]{restatabletheorem}

\title{Robust Cooperation in the Prisoner's Dilemma: //
Program Equilibrium via Provability Logic}
\author{Colomban Duclaux}
\date{\today}

\usepackage[backend=biber,style=alphabetic]{biblatex} % Ensure style matches your preference
\addbibresource{../references.bib} % Link your .bib file

\newcommand{\dupoc}{\texttt{DUPOC}}
\newcommand{\cupod}{\texttt{CUPOD}}
\newcommand{\cimcic}{\texttt{CIMCIC}}
\newcommand{\dimcid}{\texttt{DIMCID}}
\begin{document}

\maketitle
These notes recap the content exposed in \cite{barasz2021robustcooperationprisonersdilemma}

\section{Abstract}
Achieve mutual cooperation in one-shot Prisoner's Dilemma with open-source agents.
\textit{Robust}: Cooperation does not need exact equality of the agent's source code.
\textit{Unexploitable}: Agents never cooperate if the opponent defects.

\section{Introduction}
Concept of agents given full knowledge of one-another should cooperate. True for "superrationality" but not in practice. 
A line of work on \texttt{CliqueBot}: "cooperate iff the opponent's source code is identical to mine" $\to$ Nash equilibrium for the expanded game: 2 players decide which code to submit to the Prisoner's Dilemma with source-code swap.
Two faulty ideas: \begin{itemize}
    \item simulating the other agent to see what they do when given one's own source code $\to$ infinite regress when the two do the same
    \item find fixed points of strategies $\to$ sometimes several or none exist
\end{itemize}
\textbf{Solution:} Use modal logic, in particular Löb's Theorem.

Open-source Prisoner's Dilemma is analogous to Newcomb's problem.

\section{Agents in Formal Logic}
\printbibliography
\end{document}